\documentclass[11pt,a4paper]{article}
\usepackage[T1]{fontenc}
\usepackage[utf8]{inputenc}
\usepackage[french]{babel}
\usepackage{lmodern}
\usepackage{microtype}
\usepackage[margin=2.2cm]{geometry}
\usepackage{booktabs}
\usepackage{amsmath}
\usepackage{enumitem}
\usepackage[most]{tcolorbox}
\usepackage[hidelinks]{hyperref}
\setlist{itemsep=2pt,topsep=3pt}
\newtcolorbox{dire}{colback=green!4,colframe=green!40!black,title=À dire,fonttitle=\bfseries\small,boxrule=0.5pt}
\newtcolorbox{piege}{colback=red!4,colframe=red!50!black,title=Point délicat,fonttitle=\bfseries\small,boxrule=0.5pt}

\title{Guide d'explication et de soutenance\\[0.3em]\large Bande L (NISAR) contre bande C (Sentinel-1) pour la détection d'une coupe rase}
\author{Mamadou Galy Diallo}
\date{Novembre 2026}

\begin{document}
\maketitle
\tableofcontents

%=====================================================================
\section{Le récit en deux minutes}

\begin{dire}
« Le radar en bande C de Sentinel-1 est aujourd'hui l'outil standard pour détecter les coupes forestières. Depuis 2025,
NISAR fournit gratuitement de la bande L sur l'Europe, et ROSE-L prendra le relais en 2028. J'ai voulu vérifier
concrètement, sur un cas réel, ce que la bande L apporte.

J'ai d'abord inventorié les données disponibles sur les Vosges, puis cherché une vraie coupe rase de l'été 2026 avec
Sentinel-2. Sur 16 candidats, la plupart étaient des dépérissements de scolytes : j'ai retenu une coupe de 2,4~ha
près du Thillot, faite début août, et une forêt témoin autour.

J'ai comparé les produits NISAR de niveau~2, générés par ISCE3, avec 13 images Sentinel-1 que j'ai traitées sous SNAP,
le tout sur une même grille de 20~m. Pour mesurer la séparation entre coupe et forêt, j'ai utilisé l'AUC avec des
intervalles de confiance par bootstrap par blocs, et j'ai vérifié qu'avant la coupe les deux zones étaient
indiscernables.

Résultat principal : en bande L, la forêt reste cohérente d'une image à l'autre, si bien que la coupe apparaît comme
une chute nette de cohérence (AUC 0,995). En bande C, la forêt est déjà décorrélée, et la coupe ne se voit presque pas
en cohérence (AUC 0,74). En rétrodiffusion, les deux bandes voient la coupe. La bande L fait mieux dans les
géométries montantes, mais à géométrie égale (en descendante) les deux sont comparables : je le dis clairement.

Limite principale : un seul site. La suite logique, que je propose pour le stage, est de répliquer sur d'autres coupes
et sur des dépérissements, et de construire un détecteur combinant les bandes L et C. »
\end{dire}

%=====================================================================
\section{Les notions à maîtriser}

\subsection{Radar}
\begin{description}
\item[Bande, longueur d'onde.] C : 5,6~cm (5,4~GHz). L : 24~cm (1,26~GHz). P : 70~cm (BIOMASS). Une onde interagit
  surtout avec des objets de taille comparable à sa longueur d'onde : la bande C interagit avec les feuilles et les
  rameaux, la bande L avec les branches et les troncs, et elle pénètre davantage le couvert.
\item[Polarisation.] Émission et réception en H ou en V.
  \begin{itemize}
  \item HH et VV (co-polarisation) : surface et double rebond tronc-sol.
  \item HV et VH (polarisation croisée) : dépolarisation par la diffusion de volume (houppier). Elle est liée à la
    biomasse.
  \item NISAR acquiert en HH et HV sur l'Europe, Sentinel-1 en VV et VH.
  \end{itemize}
\item[$\sigma^0$, $\beta^0$, $\gamma^0$.] Coefficients de rétrodiffusion qui diffèrent par la surface de référence.
  $\gamma^0$ est normalisé par la surface perpendiculaire au faisceau ; c'est le moins dépendant de l'angle
  d'incidence.
\item[RTC / Terrain-Flattening.] Correction radiométrique du relief (Small, 2011) : en montagne, un versant face au
  radar paraît plus brillant. On normalise par la vraie surface éclairée, calculée avec un MNT. NISAR GCOV est déjà
  corrigé ; sous SNAP, c'est l'opérateur Terrain-Flattening.
\item[SLC.] \emph{Single Look Complex} : image complexe (amplitude et phase) en géométrie radar. Elle est
  indispensable pour calculer une cohérence.
\item[IW, sous-fauchée, burst.] Sentinel-1 en mode IW image 3 sous-fauchées (IW1--3), chacune découpée en bursts
  d'environ 20~km le long de la trace (technique TOPS).
\item[ESD.] \emph{Enhanced Spectral Diversity} : correction fine du recalage entre deux images TOPS, estimée dans les
  zones de recouvrement entre bursts. Avec un seul burst, il n'y a pas de recouvrement : l'ESD échoue et donne une
  cohérence nulle. Je l'ai donc retirée, ce qui est acceptable pour un site loin des jonctions de bursts.
\item[Multivue (multilook).] Moyenne de pixels voisins pour réduire le chatoiement (speckle), au prix de la
  résolution.
\end{description}

\subsection{Interférométrie}
\begin{description}
\item[Cohérence.] $\gamma=|\langle s_1 s_2^*\rangle|/\sqrt{\langle|s_1|^2\rangle\langle|s_2|^2\rangle}$, calculée sur
  une petite fenêtre. Elle vaut 1 si les diffuseurs n'ont pas bougé entre les deux dates, et tend vers 0 s'ils ont
  changé.
\item[Décorrélation temporelle.] Perte de cohérence due aux changements de la scène entre les deux dates (vent dans
  les feuilles, pluie, croissance, coupe). Elle croît avec l'écart temporel et décroît avec la longueur d'onde : la
  bande L la subit moins (Zebker et Villasenor, 1992).
\item[Biais de l'estimateur.] Sur peu de pixels, la cohérence estimée est biaisée vers le haut. Même une scène
  totalement décorrélée donne environ $\sqrt{\pi/(4N)}$ pour $N$ vues indépendantes, soit environ 0,15--0,2 avec la
  fenêtre Sentinel-1. La forêt en bande C (0,28) est donc proche de ce plancher : la coupe ne peut presque pas la
  faire baisser davantage.
\item[GUNW, GCOV.] Produits NISAR de niveau~2 générés par ISCE3 :
  \begin{itemize}
  \item GUNW : interférogramme géocodé, avec la cohérence et la phase déroulée. Il n'existe qu'en HH.
  \item GCOV : matrice de covariance géocodée, c'est-à-dire $\gamma^0$ HH et HV corrigés du relief.
  \end{itemize}
\item[Paire de 24~jours.] Fin juillet, une acquisition NISAR manque : la paire qui encadre la coupe fait 24~jours
  au lieu de 12. Plus l'écart est long, plus la décorrélation est forte, ce qui \emph{défavorise} la bande L.
\end{description}

\subsection{Statistiques}
\begin{description}
\item[AUC.] Probabilité qu'un pixel de coupe tiré au hasard ait une valeur « plus coupe » qu'un pixel de forêt tiré
  au hasard. 0,5 = hasard, 1 = séparation parfaite. Elle ne dépend d'aucun seuil ni de la forme des distributions.
  J'utilise $\max(\mathrm{AUC},1-\mathrm{AUC})$ pour ne pas dépendre du sens de la variation.
\item[d de Cohen.] Écart des moyennes divisé par l'écart-type commun. $|d|>0{,}8$ est considéré comme un « grand »
  effet ; ici on observe $|d|$ de 3 à 9.
\item[Jeffries-Matusita.] Distance entre deux distributions gaussiennes, de 0 à 2. Elle est classique en
  télédétection pour la séparabilité de classes.
\item[Pourquoi pas de p-valeurs ?] Les pixels voisins se ressemblent (autocorrélation spatiale, chatoiement, multivue) :
  ils ne sont pas indépendants. Un test classique croirait avoir 600 observations indépendantes et donnerait des
  p-valeurs absurdement petites.
\item[Bootstrap par blocs.] On rééchantillonne des blocs de $3\times3$ pixels (60~m) plutôt que des pixels isolés, ce
  qui respecte la corrélation locale (Künsch, 1989). On répète 300 fois et on prend les percentiles 2,5 et 97,5~\% :
  c'est l'intervalle de confiance.
\item[Otsu, F1, prévalence.]
  \begin{itemize}
  \item Otsu : seuil automatique qui sépare au mieux un histogramme en deux.
  \item F1 : moyenne harmonique de la précision et du rappel.
  \item Prévalence : part de la classe recherchée, ici 0,4~\% de pixels coupés. Quand elle est très faible, même un
    bon classifieur (AUC 0,9) fait plus de fausses alarmes que de vraies détections, et le F1 reste bas.
  \end{itemize}
\end{description}

%=====================================================================
\section{La méthode, choix par choix (et pourquoi)}

\begin{enumerate}
\item \textbf{Hypothèses écrites avant les calculs} (H1--H4). Cela évite de « raconter » après coup ce qu'on a
  trouvé.
\item \textbf{Site trouvé par l'optique, contrôlé à l'œil.} La référence (Sentinel-2) est indépendante du radar :
  on ne juge pas le radar avec lui-même. Le contrôle visuel a éliminé les dépérissements, qu'un simple seuil NDVI
  confondait avec des coupes.
\item \textbf{Forêt témoin locale} (anneau de 100--275~m) : même versant, même altitude et même peuplement. Elle
  absorbe les effets communs (météo, saison, humidité) : on compare la coupe à ce qu'elle serait restée.
\item \textbf{Contrôle avant coupe.} Avant la coupe, l'AUC entre coupe et témoin vaut 0,51--0,62 : les deux zones
  étaient équivalentes. Ce que l'on mesure ensuite vient bien de la coupe.
\item \textbf{Mêmes grandeurs physiques des deux côtés} : $\gamma^0$ corrigé du relief et cohérence de paires
  courtes, sur la même grille de 20~m.
\item \textbf{Sensibilité} : en enlevant les pixels de bord (érosion de 20~m), toutes les AUC augmentent. Les
  résultats ne viennent donc pas des pixels mixtes.
\item \textbf{Ingénierie} : lecture à distance des seuls blocs utiles (NISAR) et extraction partielle des zips
  (Sentinel-1) ; code en modules réutilisables, tests, tables générées automatiquement.
\end{enumerate}

%=====================================================================
\section{Lire les résultats}

\begin{itemize}
\item \textbf{Séries temporelles NISAR.} Le HV de la coupe (rouge) suit la forêt (bleu) jusqu'à fin juillet, puis
  chute de 4~dB et reste bas. La cohérence de la paire qui encadre la coupe tombe à 0,29 alors que la forêt reste
  à 0,69. Le creux commun du 28/09 au 10/10 (0,3 pour les deux zones) est un effet météo : c'est pourquoi on compare
  toujours à un témoin de la même paire.
\item \textbf{Séries Sentinel-1.} La cohérence est basse partout (0,2--0,35) : on ne voit rien de net. Le $\gamma^0$
  baisse de 2 à 4~dB en août et septembre, puis revient au niveau de la forêt en octobre.
\item \textbf{Persistance.} AUC date par date. En L montante, le contraste reste au-dessus de 0,82 jusqu'en
  novembre. En C, il s'effondre en octobre. Mon explication, plausible mais non vérifiée : à l'automne, le sol nu
  humide et les rémanents rediffusent fortement en bande C.
\item \textbf{Synthèse à géométrie égale.} C'est la figure la plus importante pour l'honnêteté du propos.
  \begin{itemize}
  \item La cohérence L gagne dans les trois géométries.
  \item En rétrodiffusion, la bande L ne gagne que dans les deux géométries montantes. En descendante (D138 contre
    D139, la comparaison la plus juste), la bande C fait un peu mieux (0,91 contre 0,84) et la persistance est
    équivalente.
  \end{itemize}
\end{itemize}

\begin{piege}
Ne pas dire « la bande L est meilleure en tout ». Dire : « l'avantage robuste de la bande L est la cohérence ; en
rétrodiffusion, il dépend de la géométrie d'acquisition, et je ne l'explique pas encore. » C'est une position plus
solide, et elle ouvre naturellement sur le travail du stage.
\end{piege}

%=====================================================================
\section{Mon regard critique sur ce travail}

\subsection{Ce qui est solide}
\begin{itemize}
\item La logique expérimentale : hypothèses a priori, référence indépendante, témoin local, contrôle avant coupe,
  sensibilité. C'est un protocole avant / après avec témoin, la bonne structure pour une étude de cas.
\item Le résultat sur la cohérence : effet énorme ($d<-3{,}7$), retrouvé dans trois géométries indépendantes,
  obtenu malgré un handicap (paire de 24~j contre 12~j) et cohérent avec la physique connue.
\item La transparence : tout est versionné, reproductible, et les chiffres du rapport viennent directement du calcul.
\end{itemize}

\subsection{Ce qu'un relecteur critiquerait}
\begin{enumerate}
\item \textbf{$n=1$ site.} C'est la limite principale. Les intervalles de confiance décrivent l'incertitude
  \emph{sur ce site}, pas la généralisation. Un second site (C09) ou une série de coupes est indispensable.
\item \textbf{Estimateurs de cohérence différents.} Le nombre de vues, la résolution et le filtrage diffèrent entre
  le GUNW NISAR et la cohérence SNAP. Le biais de l'estimateur près de zéro joue contre la bande C. Ce n'est pas une
  erreur, mais il faudrait harmoniser (même nombre de vues effectives) pour isoler l'effet de la seule longueur
  d'onde.
\item \textbf{Comparaison déséquilibrée en géométries} : trois géométries NISAR contre une seule orbite
  Sentinel-1. Il faudrait ajouter les orbites S1 montantes (15 et 88) pour comparer montante contre montante.
\item \textbf{Heures d'acquisition.} NISAR passe vers 6~h~19 (montante) et 20~h~42 (descendante), Sentinel-1 vers
  7~h~42. Rosée et vent diffèrent : cela peut expliquer une partie des différences entre géométries.
\item \textbf{Peu de dates avant la coupe} (1 à 3 paires) : la variabilité naturelle « avant » est mal estimée.
\item \textbf{Explications non vérifiées} : la remontée automnale en bande C (humidité, rémanents) et la faiblesse
  de D138 (orientation du versant par rapport à la visée) restent des hypothèses. On pourrait les tester avec la
  météo (ERA5 ou Météo-France) et la carte d'angle d'incidence local.
\item \textbf{Choix de variables multiples} : en testant beaucoup de variables, la « meilleure » est un peu
  optimiste. C'est pourquoi le message s'appuie sur les variables prévues dans les hypothèses, pas sur le maximum
  global.
\item \textbf{Produits NISAR provisoires} (2026) : la calibration peut encore évoluer.
\end{enumerate}

\subsection{Mon avis global}
Pour un test de faisabilité, la démarche est rigoureuse et la conclusion principale (la cohérence en bande L capte
l'événement de coupe là où la bande C échoue) est crédible et physiquement fondée. La conclusion sur la
rétrodiffusion doit rester prudente. Ce travail est un bon point de départ pour le stage : il démontre la maîtrise
de la chaîne (accès aux données, SNAP, ISCE3, statistiques) et identifie précisément les questions ouvertes.

%=====================================================================
\section{Questions probables et réponses}

\begin{description}
\item[Pourquoi la cohérence C est-elle si basse en forêt ?] La bande C interagit avec les feuilles et les rameaux,
  qui bougent avec le vent : en 12~jours, la forêt se décorrèle presque entièrement. La valeur mesurée (0,28) est
  proche du plancher de l'estimateur. La coupe n'a donc « rien à faire baisser ».
\item[Votre paire NISAR fait 24 jours, n'est-ce pas injuste ?] Si, mais au détriment de la bande L : une paire plus
  longue décorrèle davantage. Malgré cela, la forêt reste à 0,69 en L, contre 0,28 en C à 12~jours. Cela renforce la
  conclusion.
\item[Pourquoi pas de p-valeurs ?] Les pixels voisins ne sont pas indépendants. J'utilise des tailles d'effet (AUC,
  d) et un bootstrap par blocs, qui respecte cette dépendance.
\item[Comment savez-vous que c'est la coupe et pas autre chose ?] Avant la coupe, coupe et témoin sont
  indiscernables (AUC 0,51--0,62). Le changement apparaît exactement dans la fenêtre datée par Sentinel-2 et
  persiste ensuite.
\item[Pourquoi D138 est-il moins bon ?] Je ne l'ai pas encore démontré. Pistes : orientation du versant par rapport
  à la visée (ombre et bord de coupe), heure d'acquisition (20~h~42), et une date de moins avant la coupe. C'est
  une question que je propose d'étudier.
\item[Le F1 est faible : la méthode marche-t-elle ?] Le classement des pixels est bon (AUC de 0,82 à 0,93). Le F1
  est bas parce que les coupes ne font que 0,4~\% des pixels et que je seuille une seule variable. Un vrai
  détecteur combine des variables, filtre les petits objets et confirme sur plusieurs dates.
\item[Pourquoi avoir retiré l'ESD ?] Le site tient dans un seul burst : sans zone de recouvrement, l'ESD ne peut pas
  être estimée et SNAP renvoie une cohérence nulle. Le recalage géométrique (Back-Geocoding avec orbites précises et
  MNT) suffit loin des jonctions de bursts.
\item[Pourquoi une grille de 20~m ?] C'est la résolution native du GCOV et de la cohérence GUNW NISAR. Sentinel-1
  est multivu à environ 14~m, puis rééchantillonné. À 80~m (autre couche GUNW), la coupe ne contient qu'environ
  4 pixels indépendants : j'ai écarté cette couche des conclusions.
\item[Et ROSE-L ?] Mission Copernicus en bande L (lancement prévu en 2028). Ma chaîne repose sur des grandeurs
  standard ($\gamma^0$ RTC, cohérence à 12~j) : il suffira d'écrire un lecteur de ses produits. ISCE3, ouvert, peut
  aussi produire des L2 harmonisés.
\item[Et la bande P (BIOMASS) ?] BIOMASS n'acquiert pas sur l'Europe (zone d'exclusion SOTR). Je l'ai documenté à
  l'inventaire. Elle serait pertinente sur des forêts tropicales à forte biomasse, où la bande L sature à son tour.
\item[Coupe ou dépérissement de scolytes ?] C'est la vraie question opérationnelle dans les Vosges. Le dépérissement
  garde les troncs debout : on s'attend à une baisse modérée du HV en bande L et à une cohérence peu affectée. Ce
  serait un objectif du stage.
\item[Qu'avez-vous fait vous-même ?] Toute la chaîne : inventaire, recherche et contrôle du site, accès optimisé aux
  données, graphes SNAP, lecture des HDF5 NISAR, statistiques, figures et rapports. J'ai utilisé un assistant de
  programmation (IA) pour accélérer l'écriture du code ; les choix méthodologiques, leur vérification et leur
  interprétation sont les miens, et je peux justifier chacun d'eux.
\end{description}

%=====================================================================
\section{Chiffres à retenir}

\begin{center}
\begin{tabular}{lcc}
\toprule
 & Bande L (NISAR) & Bande C (Sentinel-1) \\
\midrule
Cohérence de la forêt à 12~j & 0,64--0,70 & 0,28 \\
Cohérence de la paire sur la coupe : coupe / forêt & 0,29 / 0,69 & 0,23 / 0,29 \\
AUC de la cohérence sur la coupe & 0,91--0,995 & 0,74 \\
$\Delta\gamma^0$ croisé, AUC (montante) & $-4{,}3$~dB, 0,97 & -- \\
$\Delta\gamma^0$ croisé, AUC (descendante) & $-2{,}3$~dB, 0,84 & $-2{,}2$~dB, 0,91 \\
AUC croisée en octobre--novembre & 0,91 (A144) / 0,76 (D138) & 0,74 \\
\bottomrule
\end{tabular}
\end{center}
Site : 2,36~ha, 59 pixels de 20~m ; témoin de 548 pixels ; coupe entre le 2 et le 11 août 2026 ; 0,4~\% de pixels
coupés dans la fenêtre de détection.

\end{document}
